% =====================================================================================================================
\section{Related work}
\label{sec:related}

\paragraph{Arrival prediction from automatic vehicle location} Bus arrival prediction from automatic vehicle location
(AVL) data has a long history. Early methods combined the last observed position with historical travel times between
stops or timepoints \parencite{lin1999experimental,cathey2003prescription}. Later work used Kalman filters
\parencite{cathey2003prescription,shalaby2004,vanajakshi2009kalman}, artificial neural networks
\parencite{chien2002,jeong2005} and, for stops served by several routes, models that share information across routes
\parencite{yu2011multiple}. \textcite{mazloumi2010variability} documented how travel time varies by time of day and
segment, the variation that segment-based methods learn. Reviews \parencite{moreiramatias2015survey,singh2022review,kumar2025review}
compare methods with each other or with the timetable. The method evaluated here belongs to the family of
segment-based historical methods \parencite{lin1999experimental,cathey2003prescription}: it learns the travel time of
each segment between consecutive stops by time band, with shrinkage towards the all-day value, and adds the position of
the vehicle.

\paragraph{Accuracy of operator real-time information} Evaluations of the predictions that operators actually show
passengers are rare. \textcite{cats2016realtime} evaluated the bus arrival information of Stockholm against AVL records
and found that errors grow with the prediction horizon and depend on the stop and the time of day. We found no study that
compares an operator's live predictions with an independent predictor on the same vehicles and moments.

\paragraph{Value of accurate information for passengers} Real-time information reduces actual and perceived waiting
times \parencite{watkins2011,dziekan2007dynamic,mishalani2006,liu2020waiting}, is valued by passengers
\parencite{ferris2010onebusaway,rahman2013users} and can increase ridership \parencite{tang2012chicago,brakewood2015nyc};
\textcite{brakewood2019} review these benefits. They depend on accuracy: wrong predictions reduce trust and lead to missed
buses \parencite{gooze2013benefits}.

\paragraph{Third-party predictors and open data} Predictors built only on public data are feasible:
\textcite{garciamaurino2024improving} improve arrival predictions using only data available through public transport
APIs, and \textcite{wessel2017retrospective} reconstruct actual passage times from a public vehicle location feed, the
kind of ground truth used here.

\paragraph{Comparing predictors} Comparing two predictors on the same targets is a classic forecasting problem. The loss
differential should be tested with errors that account for dependence \parencite{diebold1995comparing}, the scoring rule
should match the decision it serves \parencite{gneiting2011point}, and evaluation must be out of sample
\parencite{tashman2000outofsample}, with no information from the test period leaking into the predictor
\parencite{kaufman2012leakage}.
